\chapter{Introduction}

\section{Contextualization }
The continuous growth in data generation, processing, and transmission has increased the demand for storage systems capable of handling large volumes of information in an efficient, durable, and energy-sustainable manner. Although magnetic, optical, and semiconductor media remain predominant, they present limitations related to storage density, energy consumption, physical space, and the infrastructure required for operation and cooling. This scenario has motivated the investigation of new storage media and computational architectures capable of exploring unconventional alternatives.

Among these alternatives, synthetic DNA data storage has been studied as a long-term archival technology. National initiatives, such as research conducted by the Institute for Technological Research (IPT) in partnership with technology companies, demonstrate the growing interest in this area. However, using DNA as a storage medium does not depend solely on the physical properties of the molecule. It requires computational systems capable of converting digital data into nucleotide sequences, introducing redundancy, organizing information for retrieval, and reconstructing the original file after sequence reading.

One of the best-known architectures for this purpose is DNA Fountain \cite{erlich2017}, which is based on LT (\emph{Luby Transform}) codes \cite{luby2002}. In this approach, a digital file is divided into segments and used to produce redundant encoded units called \emph{droplets}. Each \emph{droplet} is obtained by combining a subset of the original file segments through the XOR operation. The resulting redundancy allows the data to be recovered from a sufficient number of valid \emph{droplets}, even when some of them are unavailable during decoding.

From a Software Engineering perspective, DNA Fountain can be analyzed as an encoding and decoding \emph{pipeline} composed of segmentation, \emph{droplet} generation, degree selection, XOR operations, and file reconstruction. Decoding is performed through the \emph{peeling decoder}, an iterative algorithm that identifies \emph{droplets} associated with a single unknown segment and uses this information to simplify the remaining encoding relations. This process can be represented by a bipartite graph between segments and \emph{droplets}, as well as by systems of equations over the binary finite field \(GF(2)\).

LT codes were originally proposed for erasure channels, in which an encoded unit is either received correctly or not received at all. However, the physical processes involved in DNA storage include stages such as synthesis, storage, amplification, and sequencing, which may introduce more diverse types of errors. Detailed modeling of these processes requires experimental knowledge, reference data, and models specific to each technology. Therefore, this work does not intend to reproduce the biological behavior of the DNA channel. Instead, it uses an abstract probabilistic erasure channel as a controlled approximation for studying the behavior of DNA Fountain.

In this model, each \emph{droplet} generated by the encoder may be retained or removed according to an experimentally defined erasure probability. This abstraction makes it possible to investigate, in a reproducible manner, how the number of available \emph{droplets}, the \emph{overhead}, and the parameters of the robust soliton distribution influence the probability of complete file recovery. This scope makes it possible to concentrate the analysis on the relationship between LT code theory, the behavior of the \emph{peeling decoder}, and the implementation of a software simulator.

In this context, this work lies at the intersection of Software Engineering, information theory, coding theory, and algorithmic bioinformatics. The concepts of sequences, graphs, and computational problem formulation presented by Compeau and Pevzner \cite{compeau2018} provide support for representing data and encoding relationships in a structured manner. However, the foundations of LT codes, the robust soliton distribution, and the \emph{peeling decoder} are addressed specifically, as these elements belong to the domain of Fountain codes and erasure channels.

\section{Motivation}

\section{Research Problem}

\section{Objectives}

The general objective of this work is to analyze, design, and evaluate a modular DNA Fountain pipeline simulator capable of representing LT codes, biological constraints, erasure losses, and peeling decoding, allowing experimental investigation of system parameters.

The specific objectives are:

\begin{enumerate}
    \item To formalize the foundations of LT codes, the robust soliton distribution, and the \emph{peeling decoder} in the context of DNA data storage.
    
    \item To computationally model the essential stages of the DNA Fountain pipeline, including segmentation, droplet generation, seed-based headers, and deterministic index reconstruction.
    
    \item To implement a modular Python simulator with a clear separation of concerns across distribution, coding, channel, biological screening, decoding, and metrics collection.
    
    \item To enforce biological constraints (such as GC-content and homopolymer limits) via a configurable screening filter.
    
    \item To simulate erasure channels and invalid read discard mechanisms.
    
    \item To experimentally evaluate the impact of overhead, erasure probability, block size, and robust soliton distribution parameters ($c, \delta$) on complete data recovery and computational cost.
    
    \item To ensure reproducibility through explicit seed control, batch experimentation, and configuration logging.
\end{enumerate}

\section{Delimitation of Scope}

Accordingly, this work is delimited as a Software Engineering investigation applied to the computational study of Fountain codes for DNA data storage. The expected contribution consists of building a reproducible simulator, analyzing DNA Fountain parameters, and evaluating recovery conditions over erasure channels. The scope does not include DNA molecule synthesis, physical sequencing, laboratory analysis of samples, biologically faithful modeling of synthesis and reading processes, or the complete implementation of a commercial molecular storage system.

As a continuation of the project, the implementation of a restricted \emph{peeling decoder} core on a Digilent Basys 3 FPGA may be investigated. This stage would aim to evaluate the acceleration of decoding operations using reconfigurable hardware, based on instances generated by the simulator developed in this work. This possible continuation does not alter the scope of TCC I, which remains focused on modeling, simulator implementation, and experimental analysis in software.

\section{Organization of the Work}

The remainder of this work is organized as follows. Chapter~\ref{chap:theoretical_foundation} presents the theoretical foundation of DNA data storage, error models, fountain codes, and decoding algorithms. Chapter~\ref{chap:related_work} reviews related work in DNA Fountain, molecular storage simulators, and identifies the existing gap. Chapter~\ref{chap:software_architecture} discusses the system requirements, architectural decisions, and modular software engineering of the developed simulator. Chapter~\ref{chap:experimental_methodology} details the experimental design, variables, and evaluation scenarios. Chapter~\ref{chap:results} presents the experimental results and discussion. Finally, Chapter~\ref{chap:conclusion} presents concluding remarks, and Chapter~\ref{chap:future_work} outlines perspectives for future work.